\chapter{Bernstein Polynomials Library}
In Figure~\ref{fig:bernstein_library_class_diagram} we can see the architecture behind the Bernstein Polynomials Library.
\begin{figure}[htbp]
    \centering
    \includegraphics[max width=\textwidth, max height=0.85\textheight, keepaspectratio]{images/2-bernstein/class_diagram.pdf}
    \caption{Bernstein Polynomials Library Class Diagram.}
    \label{fig:bernstein_library_class_diagram}
\end{figure}

%%%%%%%%%%%%%%%%%%%%%%%%%%%%%%%%%%%%%%%%%%%%%%%%%
\section{Library Constraints}
\label{sec:bernstein_library_constraints}
The critical part of the library, from a computational perspective, is the calculation of the binomial coefficients. The whole library is based on \texttt{double} type and this implies some limitations on the degree of the polynomial that can be evaluated.

\myskip

These limitations are due to two factors:
\begin{itemize}
    \item Range \\
        The exponent allows us to represent numbers up to $1.8\times10^{308}$.
    \item Precision \\
        We can guarantee a precision up to 15 decimal digits: the 52 digits of the mantissa brings the machine $\epsilon$ error to circa $10^{-16}$.
\end{itemize}

\myskip

The consequence of this approach is the following:
\begin{itemize}
    \item Within a $degree<54$ we can compute the polynomial without any precision loss.
    \item If the degree is within the $[54,1029]$ interval we have a precision loss due to the large binomial coefficients and small powers, but we can still compute the polynomial.
    \item For $degree\ge 1030$ we have an overflow of the binomial coefficients and the polynomial cannot be computed.
\end{itemize}

%%%%%%%%%%%%%%%%%%%%%%%%%%%%%%%%%%%%%%%%%%%%%%%%%
\section{Approximation of Empirical Cumulative Distribution Function}
In this Section we discuss the part of the library responsible for the approximation of a Empirical Cumulative Distribution Function (ECDF) with a Bernstein polynomial.

The system allows us to approximate both CDF and PDF curve, given an empirical CDF $F$. For each curve we have three variants of basis:
\begin{itemize}
    \item \textbf{Normal}: The original ECDF must be in $[0,1]$, i.e. $F:[0,1]\to[0,1]$.
    \item \textbf{Linear}: The original ECDF must be in $[a,b]$, i.e. $F:[a,b]\to[0,1]$, where both $a$ and $b$ must be finite value.
    \item \textbf{Exponential}: The original ECDF can be in $[0,+\infty]$, i.e. $F:[0,+\infty]\to[0,1]$.
\end{itemize}

\subsection{Implementation}
From the Figure~\ref{fig:bernstein_library_class_diagram} we can see that there are only two single classes that work together:
\begin{itemize}
    \item \textbf{BernsteinFromECDF} \\
        It is the facade of the system. It exposes six methods, three for the approximation of the empirical CDF and three for the empirical PDF.
    \item \textbf{ECDF} \\
        It is the representation of the empirical CDF. It allows the evaluation of the ECDF in a point $x$.
\end{itemize}

%%%%%%%%%%%%%%%%%%%%%%%%%%%%%%%%%%%%%%%%%%%%%%%%%
\section{Approximation of Polynomial}
The system allows us to map a polynomial in monomial form to a Bernstein polynomial. In particular, given a polynomial $\displaystyle P(x) = \sum_{i=0}^{M} a_i x^i$ of degree $M$ with known analytical form, we want to identify the Bernstein Polynomial $\displaystyle B(x) = \sum_{i=0}^{N} \beta_i b_{i,N}(x)$ of degree $N$ that best fits it.

\subsection{Implementation}
In Figure~\ref{fig:bernstein_library_class_diagram} there is the subsystem that is responsible for the approximation from the polynomial defined by monomial to a Bernstein Polynomial. Let's see the main components:
\begin{itemize}
    \item \textbf{BernsteinFromMonomial} \\
        It's the facade to use the system and abstract the domain model that shape the polynomial.
    \item \textbf{Polynomial} \\
        It's the abstract class that wraps the state of a polynomial, i.e. the coefficients that define the polynomial. It defines also the \texttt{eval} method to evaluate the polynomial in given point \texttt{x} and the method \texttt{convolve()} to perform the convolution between two polynomials of the same type.
    \item \textbf{MonomialPolynomial} \\
        Classical polynomial where the basis is defined as an integer power of \texttt{x}.
    \item \textbf{BernsteinPolynomial} \\
        It represents the Bernstein polynomial. It's defined by the coefficient and the type of \texttt{Basis}. There is only one kind of basis: the linear basis map the input \texttt{x} in the \texttt{[a,b]} interval.
\end{itemize}

\subsection{Approximation Techniques}
We have tested some mapping techniques from the polynomial in monomial form and the Bernstein polynomial.

\subsubsection{Support in $[0,1]$ and $N=M$}
\label{subsec:monomial_approximation_support_0_1_N_equal_M}
The simplest scenario is if we have a support in $[0,1]$ and the degree of the input and output polynomials are the same.
\begin{itemize}
    \item We can provide a closed form formula to perform the conversion of the coefficients from the monomial ones to the Bernstein ones. This formula is $\displaystyle\beta_i = \sum_{j=0}^{i} \frac{\binom{i}{j}}{\binom{N}{j}} a_j$. We have a $\mathcal{O}(n^2)$ cost due to two inner loops.
    \item We can multiply the vector of Bernstein coefficients $b$ and a conversion matrix $M$ to obtain the monomial coefficient $a$ as $a=Mb$; in this way we can find the Bernstein coefficient as $b=M^{-1}a$. We have a $\mathcal{O}(n^3)+\mathcal{O}(n^2)$ cost due to the inversion of $M$ and two inner loops for the product; the inversion of the matrix is unstable for big order of $N$. This basis change matrix is obtained from these steps: the Bernstein basis are defined as $\displaystyle B_{i,N}(x) = \binom{N}{i} x^i (1-x)^{N-i}$; we can expand this in $\displaystyle B_{i,N}(x) = \binom{N}{i} x^i \sum_{j=0}^{N-i} \binom{N-i}{j} (-1)^j x^j$ and then obtain $\displaystyle B_{i,N}(x) = \binom{N}{i} x^i \left[ \sum_{j=0}^{N-i} \binom{N-i}{j} (-1)^j x^j \right]$; we can now change $i = k + j$ and obtain $\displaystyle B_{i,N}(x) = \sum_{j=i}^{N} \left[ \binom{N}{i} \binom{n-i}{j-i} (-1)^{j-i} \right] x^j$; at this point we can define the matrix as $M_{i,j} = \begin{cases} \binom{n}{j} \binom{n-j}{i-j} (-1)^{i-j} & \text{ se } i \ge j \\ 0 & \text{ se } i < j \end{cases}$.
\end{itemize}

\subsection{Functionality}
This part of the library allows some operations that are independent from the domain model.
\subsection{Distance}
In the system there's the possibility of testing the distance between two given polynomials and an interval $[a,b]$. This is useful if we test the polynomial $P_m$ in monomial form and its Bernstein approximation $P_b$.

The distances implemented are the three most used types of norm, that consider the whole interval. We have considered the following:
\begin{itemize}
    \item $L_1$: It measures the global mean error and it's obtained as $\displaystyle\int_a^b |P_m(x) - P_b(x)| dx$.
    \item $L_2$: It penalizes the big error in the whole interval $[a,b]$ and it's obtained as $\sqrt{\displaystyle\int_a^b (P_m(x) - P_b(x))^2 dx}$.
    \item $L_\infty$: It represents the biggest error in the interval between the two polynomials and it's obtained as $\max_{x \in [a,b]} |P_m(x) - P_b(x)|$.
\end{itemize}

\myskip

Each of these norms, except for $L_\infty$, is calculated with a numerical integration: it is a family of techniques to approximate the value of a definite integral without calculating the primitive of the function. This is useful when it is difficult to derive the primitive of the function or when the function isn't known in the analytical form.

In this case we used a specific Newton-Cotes formula, the Simpson's rule.

\subsubsection{Visualization}
There is an abstract class that allows us to plot a polynomial with a given precision and in a given interval.